% Pursuing Stacks (cotes 134-2 to 134-8): archive cote 134-3, site page 39 (PDF page 40), typescript page 39.
% typescript. Draft compared: ps2.tex:1894-1966.
% AI-transcribed by Claude Opus 5.5 (claude-opus-5-5, adaptive thinking, effort high), Message Batches API, batch ps1 (msgbatch_01Rv5hCaNTtUt4JHx1yqoHPL), 2026-11-12; not yet checked by eye.
\scanpage{134-3/39}
\margin{\unsure{cf théorème p 46}}
\mspage{\hand{39}} \margin{\unsure{K}}

Here are the conditions I got:
(T 1) A is aspheric .
(T 2\over{)}{\unsure{)}} For $a,b \in$ Ob A, $A_{/axb}$ is aspheric \inserted{\unsure{i.e.} for any $a \in \mathrm{Ob}\, A$, $A_{/a} \to A$ is aspheric} (NB axb need not exist in A
but it is in any case well defined as an element of $\overset{\hand{\wedge}}{A}$ ).
(T 3) There exists an \uline{aspheric} element I of $\overset{\hand{\wedge}}{A}$ , and two subobjects
\margin{NB \unsure{In view of} \unsure{T1 $\to$ T2}, I aspheric $\Leftrightarrow$ $I \to e$ \unsure{is} aspheric}
$e_o$ and $e_1$ of I which are final elements of $\hat{A}$ , such that
$e_o \hand{\cap} e_1$ $(\overset{def}{=} e_o \hand{\times}_I e_1) = \hand{\emptyset}_{A^\wedge}$ , the initial or ``empty'' element of $A\hand{\hat{}}$ .

In case when A = $\hand{\Delta}$ (as we\hand{l}l as in the cubical analogon $\hand{\Box}$ ), I took
I = $\hand{\Delta}_1$ which is an element of A itself, and \deleted{the} moreover A has a final
element (which is a final element of $A^\wedge$\hand{/}therefore) e , thus $e_o$ and $e_1$ are
defined by well defined arrows in A itself, namely $\hand{\delta}_o$ and $\hand{\delta}_1$ . But it
does not seem that these special features are really relevant. In any case
\uline{intuitively I stands for the unit interval}, \uline{with endpoints} $e_o$ , $e_1$ . If
F is any element in $A^\wedge$, the standard way for trying to prove it is asphe-%
ric would be to prove that we can find a ``constant map'' F $\hand{\to}$ F, namely one
\inserted{which} factors into
\margin{\unsure{$\mathcal{A}$}}
\[ F \hand{\longrightarrow} e \overset{c}{\hand{\longrightarrow}} F \qquad \text{(e the final element of \deleted{\illegible{1}} \hand{$A^\wedge = \hat{A}$)}} \]
for suitable \hand{c} : e $\hand{\to}$ F or ``section'' of F, which be ``homotopic'' to the
identity map F$\hand{\to}$F. When trying to make explicit the notion of a ``homotopy''
\inserted{h} between two such maps, more generally between two \inserted{$f_o, f_1$ :} maps F $\rightrightarrows$ G , we
hit of course upon the arrow h in the following diagram, which should
make it commutative

\hand{(D)}
\diagram{
\begin{tikzcd}
 & F \times I \ar[dd, dashed, "\hand{h}"'] & \\
F \simeq F \times e_o \ar[ur] \ar[dr, "f_o"'] & & F \simeq F \times e_1 \ar[ul, bend right] \ar[dl, "f_1"] \\
 & G &
\end{tikzcd}
}
.

This notion of a homotopy is defined in any category \inserted{$\mathcal{A}$} where we got an
element I and two subobjects $e_o$ , $e_1$ which are final objects. Suppose we
\inserted[typed]{got such h, and we} know moreover \xout{that} (for given F,G,$f_o$,$f_1$,h) that the two inclusions of
$Fxe_o$ and $Fxe_1$ into FxI are weak equivalences, and that $f_o$ is a weak equi-%
valence (for a given set of arrows called ``weak equivalences'', for instan-%
ce defined in terms of a``topology'' on \hand{$\mathcal{A}$} , in the present case the canoni-%
cal topology of the topos $A\hand{\hat{}}$ ) ; then it follows (with the usual ``mild
saturation\unsure{\deleted{''}} condition'' on the notion of weak equivalence) that h, and hence
$f_1$ are weak equivalences. Coming back to the case F=G, $f_o$ = $id_F$ , $f_1$ =
``constant map'' defined by a c: e $\hand{\to}$ F , we get that this constant map
$f_1$ is a weak equivalence. Does this imply that F $\hand{\to}$ e is equally a weak
equivalence ? This is not quite formal for general (\hand{$\mathcal{A}$}, W), but it is \xout{\illegible{1}}
true though in the case \hand{$\mathcal{A}$} = $A^\wedge$ and with the usual meaning of ``weak equiva-%
